\documentclass[acmengage]{acmart}
\settopmatter{printacmref=false, printccs=false, printfolios=true}
\setcopyright{none}
\renewcommand\footnotetextcopyrightpermission[1]{}
\renewcommand{\shortauthors}{}
\renewcommand\acmBooktitle{}
\pagestyle{plain}

\usepackage{listings}
\usepackage{xcolor}
\usepackage{tcolorbox}

\AtBeginDocument{%
  \providecommand\BibTeX{{%
    Bib\TeX}}}

\begin{document}

\fancyhead{} % Clear all headers
\fancyfoot{} % Clear all footers
\fancyfoot[C]{\thepage} % Keep page numbers centered at bottom


\title{Latency-Fairness vs. Throughput Tradeoffs in Non-Preemptive Shared GPU Scheduling}

\author{Jianing Ou}
\email{jou32@wisc.edu}
\affiliation{%
  \institution{University of Wisconsin-Madison}}

\author{Xuechun Jin}
\email{xjin252@wisc.edu}
\affiliation{%
  \institution{University of Wisconsin-Madison}}

\author{Yayen Lin}
\email{lin383@wisc.edu}
\affiliation{%
  \institution{University of Wisconsin-Madison}}

\author{Matthew Kearney}
\email{mtkearney@wisc.edu}
\affiliation{%
  \institution{University of Wisconsin-Madison}}



\begin{abstract}
Modern cloud GPU platforms increasingly share hardware across multiple concurrent workloads, from latency-sensitive inference to throughput-oriented training. Because GPUs are fundamentally non-preemptive (a kernel cannot be interrupted once it starts), scheduling decisions critically impact job waiting times, fairness, and throughput. Understanding these tradeoffs is essential for designing effective shared GPU systems.

We study the tradeoffs among latency, fairness, and throughput in non-preemptive GPU scheduling through comprehensive empirical evaluation. We use ISCAS-85 circuit benchmarks to generate realistic task graphs with diverse dependency structures, mapping each circuit gate to a synthetic CUDA kernel that emulates specific computational characteristics (compute-bound, memory-bound, latency-sensitive) while preserving realistic dependencies. We implement and evaluate six scheduling policies: FIFO, priority-based, and four dependency-aware variants (HighFanout, CriticalPath, LevelAware, Hybrid) across sequential and concurrent execution modes. Using real GPU hardware, we measure latency, throughput, GPU utilization, and fairness (Jain's index) across balanced and imbalanced workloads with varying batch sizes (32, 128, 512).

Our experiments reveal several key findings. In balanced workloads (similar-sized circuits competing), all schedulers achieve near-perfect fairness with Jain's index above 0.96, but imbalanced workloads (mixing small and large circuits) expose fundamental tradeoffs where fairness drops to around 0.65-0.69 regardless of policy. The concurrent execution mode dramatically improves GPU utilization from 9-11\% (sequential) to 65-100\% (concurrent) while also increasing throughput, though at the cost of higher variance in task completion times.
\end{abstract}

\keywords{GPU scheduling, non-preemptive scheduling, fairness, throughput, latency, task graphs, computer architecture}

\maketitle
\section{Introduction}

\section{Background}
\label{sec:background}

\section{Related Work}
\label{sec:relatedwork}

\section{Methodology}
\label{sec:methodology}

\section{Evaluation}
\label{sec:evaluation}

\section{Conclusion}
\label{sec:conclusion}


%% the bibliography file
\bibliographystyle{ACM-Reference-Format}
\bibliography{base}

\end{document}
